\chapter{Distributions and Sobolev Spaces}

\begin{center}\itshape
How do you differentiate a function with a corner, a jump, or no derivative anywhere? You stop asking for its values and ask how it
acts. Every locally integrable function, every measure, and some stranger objects become infinitely differentiable, and
differential equations acquire solutions that classical calculus could not even write down.
\end{center}

Physicists did not wait for permission. Oliver Heaviside, a self-taught telegraph engineer, solved the equations of electrical
circuits in the 1890s with an ``operational calculus'' in which the derivative of a unit step was an infinitely tall spike. The Royal
Society published the first two papers of his series on operational methods and refused the third (1894). In 1927 Paul Dirac needed the same spike in quantum mechanics, wrote it $\delta(x)$, and stated
its rules: $\delta(x)=0$ for $x\ne0$ and $\int\delta(x)f(x)\dd x=f(0)$. No function can do this, and Dirac knew it; he called $\delta$
``improper'' and used it anyway, to great effect. In the 1930s Sergei Sobolev in Leningrad, working on wave equations, defined
\emph{generalised solutions} by integration against smooth test functions, and built the spaces that now carry his name.

The unifying idea came to Laurent Schwartz, a young French mathematician who had spent the war hiding under a false name. One night in
November 1944 he understood that all these objects are continuous linear functionals on a space of smooth functions with compact
support, and that differentiation can be moved onto the test function by integration by parts. He later recalled it as one of the most
beautiful nights of his life. His \emph{Th\'eorie des distributions} (1950--1951) earned him the Fields Medal in 1950. Distributions
gave the Fourier transform its natural domain, gave $\delta$ a precise meaning, and gave partial differential equations a framework in
which existence can be proved by Hilbert space methods and smoothness established afterwards.

This chapter builds distributions on open sets of $\R^d$, their derivatives and Fourier transforms, finds the fundamental solution of
the Laplacian, and then studies Sobolev spaces on the line and on intervals in full: embedding, Rellich's compactness, and the
Lax--Milgram theorem, which solves boundary value problems in their weak form and leads straight to the finite element method.

\section{Test functions and distributions}

Let $\Omega\subset\R^d$ be open. A \emph{test function} is a $C^\infty$ function on $\Omega$ with compact support in $\Omega$; they form the space $\mathcal D(\Omega)=C_c^\infty(\Omega)$. Mollifiers (Tome~6, Chapter~9) show that there are
many: for every compact $K\subset\Omega$ there is $\chi\in\mathcal D(\Omega)$ with $0\le\chi\le1$ and $\chi=1$ on a neighbourhood of $K$ (mollify the indicator of a small neighbourhood of $K$). For $\varphi$ smooth and $N\in\N$, write
$\|\varphi\|_{C^N}=\max_{|\alpha|\le N}\sup|\partial^\alpha\varphi|$.

\begin{definition}
A \emph{distribution} on $\Omega$ is a linear form $T:\mathcal D(\Omega)\to\C$, written $\varphi\mapsto\langle T,\varphi\rangle$, such that for every compact $K\subset\Omega$ there are $C\ge0$ and $N\in\N$ with
\[|\langle T,\varphi\rangle|\le C\|\varphi\|_{C^N}\qquad\text{for all }\varphi\in\mathcal D(\Omega)\text{ supported in }K.\]
If one $N$ works for all $K$, the least such $N$ is the \emph{order} of $T$. The distributions form a vector space $\mathcal D'(\Omega)$. A sequence $T_n\to T$ in $\mathcal D'(\Omega)$ if $\langle T_n,\varphi\rangle\to\langle T,\varphi\rangle$ for every
test function $\varphi$.
\end{definition}

The estimate is a continuity condition: if test functions $\varphi_n$ are supported in a fixed compact set and $\varphi_n\to\varphi$ with all derivatives uniformly, then $\langle T,\varphi_n\rangle\to\langle T,\varphi\rangle$.

\begin{example}\label{ex:distexamples}
(a) A locally integrable $f\in L^1_{\mathrm{loc}}(\Omega)$ defines $\langle T_f,\varphi\rangle=\int f\varphi$, a distribution of order $0$: $|\langle T_f,\varphi\rangle|\le\bigl(\int_K|f|\bigr)\|\varphi\|_{C^0}$. We write $f$ for $T_f$.
(b) A Radon measure $\mu$ (finite on compact sets) defines $\langle\mu,\varphi\rangle=\int\varphi\dd\mu$, also of order $0$. In particular the \emph{Dirac mass} $\langle\delta_a,\varphi\rangle=\varphi(a)$; we write $\delta=\delta_0$.
(c) On $\R$, the \emph{principal value}
\[\Bigl\langle\operatorname{p.v.}\frac1x,\varphi\Bigr\rangle=\lim_{\varepsilon\to0}\int_{|x|>\varepsilon}\frac{\varphi(x)}x\dd x=\int_0^\infty\frac{\varphi(x)-\varphi(-x)}x\dd x\]
is a distribution of order at most $1$: if $\varphi$ is supported in $[-R,R]$, $|\varphi(x)-\varphi(-x)|\le2|x|\|\varphi'\|_\infty$, so the integral is at most $2R\|\varphi\|_{C^1}$. It is not of order $0$
(Exercise~\ref{ex:pvorder}).
\end{example}

\begin{proposition}\label{prop:lloc}
If $f\in L^1_{\mathrm{loc}}(\Omega)$ and $\int f\varphi=0$ for every $\varphi\in\mathcal D(\Omega)$, then $f=0$ almost everywhere. So $f\mapsto T_f$ is injective, and $L^1_{\mathrm{loc}}(\Omega)\subset\mathcal D'(\Omega)$.
\end{proposition}
\begin{proof}
Let $B$ be a closed ball in $\Omega$ and $\delta>0$ with the $\delta$-neighbourhood of $B$ in $\Omega$. Let $g=f\mathbf 1_{B'}$, where $B'$ is the closed $\delta$-neighbourhood of $B$; $g\in L^1(\R^d)$. For a mollifier $\varphi$ and
$t<\delta$, $g*\varphi_t(x)=\int f(y)\varphi_t(x-y)\dd y=0$ for $x\in B$, since $y\mapsto\varphi_t(x-y)$ is a test function supported in $B'$. As $g*\varphi_t\to g$ in $L^1$ (Tome~6, Theorem~9.10), $g=0$ almost everywhere on $B$.
Countably many balls cover $\Omega$.
\end{proof}

\begin{think}
Why is the Dirac mass not a function? Suppose $\int f\varphi=\varphi(0)$ for all $\varphi$, with $f$ locally integrable. Test it on functions supported away from $0$, and use Proposition~\ref{prop:lloc}.
\end{think}

(Answer: $f=0$ almost everywhere on $\R^d\setminus\{0\}$, hence everywhere almost, and then $\varphi(0)=0$ for every $\varphi$, which is absurd.)

\section{Derivatives}

If $f\in C^1(\Omega)$ and $\varphi\in\mathcal D(\Omega)$, integration by parts in the variable $x_j$ (with no boundary terms, since $\varphi$ has compact support) gives $\int\partial_jf\,\varphi=-\int f\,\partial_j\varphi$. This formula makes
sense for every distribution, and becomes the definition.

\begin{definition}
For $T\in\mathcal D'(\Omega)$ and a multi-index $\alpha$, $\langle\partial^\alpha T,\varphi\rangle=(-1)^{|\alpha|}\langle T,\partial^\alpha\varphi\rangle$. For $a\in C^\infty(\Omega)$, $\langle aT,\varphi\rangle=\langle T,a\varphi\rangle$.
\end{definition}

Both are distributions: $\|\partial^\alpha\varphi\|_{C^N}\le\|\varphi\|_{C^{N+|\alpha|}}$, and by Leibniz $\|a\varphi\|_{C^N}\le C_{K,N}\|\varphi\|_{C^N}$ for $\varphi$ supported in $K$. So \emph{every distribution is infinitely differentiable},
derivatives commute ($\partial_i\partial_j\varphi=\partial_j\partial_i\varphi$), and the Leibniz rule $\partial_j(aT)=(\partial_ja)T+a\,\partial_jT$ holds (test it on $\varphi$). For $f\in C^1$, the distributional and classical derivatives
agree.

\begin{example}\label{ex:derivatives}
On $\R$, let $H=\mathbf 1_{[0,\infty[}$ be the Heaviside function. Then $\langle H',\varphi\rangle=-\int_0^\infty\varphi'=\varphi(0)$: $H'=\delta$. Similarly $\langle\delta',\varphi\rangle=-\varphi'(0)$. The function $\log|x|$ is locally
integrable, and
\[\langle(\log|x|)',\varphi\rangle=-\lim_{\varepsilon\to0}\int_{|x|>\varepsilon}\log|x|\,\varphi'(x)\dd x=\lim_{\varepsilon\to0}\Bigl(\int_{|x|>\varepsilon}\frac{\varphi(x)}x\dd x+\log\varepsilon\,\bigl(\varphi(\varepsilon)-\varphi(-\varepsilon)\bigr)\Bigr),\]
and the last term tends to $0$ since $|\varphi(\varepsilon)-\varphi(-\varepsilon)|\le2\varepsilon\|\varphi'\|_\infty$. So $(\log|x|)'=\operatorname{p.v.}\frac1x$.
\end{example}

More generally, a function $f$ which is $C^1$ on $\R$ except for a jump at $a$, with $f'$ (where it exists) locally integrable, has distributional derivative $f'+(f(a^+)-f(a^-))\delta_a$: the jump becomes a Dirac mass.

\begin{theorem}\label{thm:zeroderiv}
Let $I$ be an open interval and $T\in\mathcal D'(I)$ with $T'=0$. Then $T$ is a constant function. If $T'=g$ for a continuous $g$, then $T=G+c$ for any primitive $G$ of $g$ and a constant $c$.
\end{theorem}
\begin{proof}
A test function $\psi$ is the derivative of a test function if and only if $\int\psi=0$: if $\int\psi=0$, the primitive $x\mapsto\int_{-\infty}^x\psi$ (with $\psi$ extended by $0$) vanishes to the left of the support of $\psi$ and, since
$\int\psi=0$, to its right. Fix $\theta\in\mathcal D(I)$ with $\int\theta=1$. For every $\varphi\in\mathcal D(I)$, $\varphi-(\int\varphi)\theta$ has integral $0$, so it is $\psi'$ with $\psi\in\mathcal D(I)$, and
\[\langle T,\varphi\rangle=\Bigl(\int\varphi\Bigr)\langle T,\theta\rangle+\langle T,\psi'\rangle=c\int\varphi-\langle T',\psi\rangle=c\int\varphi,\]
with $c=\langle T,\theta\rangle$. So $T=c$. For the second claim, $(T-G)'=g-g=0$, as the classical derivative of the $C^1$ function $G$ is also its distributional derivative.
\end{proof}

\section{Smoothing}

For $T\in\mathcal D'(\R^d)$ and $\varphi\in\mathcal D(\R^d)$, define the \emph{convolution} $(T*\varphi)(x)=\langle T,\varphi(x-\cdot)\rangle$. For $T=T_f$ it is the usual $f*\varphi$.

\begin{proposition}\label{prop:smoothing}
$T*\varphi$ is $C^\infty$, with $\partial^\alpha(T*\varphi)=T*\partial^\alpha\varphi=(\partial^\alpha T)*\varphi$. If $\varphi$ is a mollifier, then $T*\varphi_t\to T$ in $\mathcal D'(\R^d)$ as $t\to0$. So every distribution is a limit of smooth
functions.
\end{proposition}
\begin{proof}
Fix $x$ and a unit vector $e_j$. As $h\to0$, the test functions $\frac1h(\varphi(x+he_j-\cdot)-\varphi(x-\cdot))$ are supported in a fixed compact set and converge to $\partial_j\varphi(x-\cdot)$ with all their derivatives, uniformly, by
Taylor's formula applied to each derivative of $\varphi$. By the continuity estimate, $\partial_j(T*\varphi)(x)$ exists and equals $(T*\partial_j\varphi)(x)$; the same argument gives continuity of $T*\varphi$ and, by
induction, all derivatives. And $\langle\partial_jT,\varphi(x-\cdot)\rangle=-\langle T,\partial_{y_j}[\varphi(x-y)]\rangle=\langle T,(\partial_j\varphi)(x-\cdot)\rangle$.

For the convergence, let $\psi\in\mathcal D$ and $\check\varphi_t(y)=\varphi_t(-y)$. A Riemann sum argument (the same continuity estimate, applied to Riemann sums of $\int\varphi_t(x-y)\psi(x)\dd x$, which converge to it in the sense of
test functions) gives $\langle T*\varphi_t,\psi\rangle=\langle T,\check\varphi_t*\psi\rangle$. As $t\to0$, $\check\varphi_t*\psi\to\psi$ with all derivatives uniformly (Tome~6, Theorem~9.10(b), applied to $\partial^\alpha\psi$), with
supports in a fixed compact set, so $\langle T,\check\varphi_t*\psi\rangle\to\langle T,\psi\rangle$.
\end{proof}

\begin{figure}[ht]
\centering
\fig{ch06-dirac}
\caption{The mollifiers $\varphi_t$ for $t=1$, $\frac12$ and $\frac14$, each of integral $1$. As $t\to0$ they converge to $\delta$ in $\mathcal D'(\R)$, although they converge pointwise to $0$ away from the origin and to
$+\infty$ at it: convergence of distributions is convergence of averages, not of values.}
\label{fig:dirac}
\end{figure}

\section{Tempered distributions and the Fourier transform}

The transform of a test function is never a test function (Corollary~\ref{cor:nocompact}), so the Fourier transform cannot act on $\mathcal D'$. It acts on the distributions that are continuous for the topology of the Schwartz space.
For $\varphi\in\mathcal S(\R^d)$ write $p_N(\varphi)=\max_{|\alpha|,|\beta|\le N}\sup_x|x^\alpha\partial^\beta\varphi(x)|$.

\begin{definition}
A \emph{tempered distribution} is a linear form $T$ on $\mathcal S(\R^d)$ with $|\langle T,\varphi\rangle|\le Cp_N(\varphi)$ for some $C,N$. They form the space $\mathcal S'(\R^d)$.
\end{definition}

Restricted to $\mathcal D$, a tempered distribution is a distribution, and it is determined by this restriction since $\mathcal D$ is dense in $\mathcal S$ for the seminorms $p_N$ (multiply by $\chi(\frac xR)$ with $\chi=1$ near $0$). Functions in
$L^p$, $1\le p\le\infty$, measurable functions with polynomial growth, finite measures, $\delta$ and $\operatorname{p.v.}\frac1x$ are tempered. The function $e^x$ is not: with $g(x)=e^{-\sqrt{1+x^2}}\in\mathcal S$ and $\chi\in\mathcal D$ equal to $1$ near $0$, the functions $\chi(\frac xn)g(x)$ have bounded
seminorms $p_N$, while $\int e^x\chi(\frac xn)g(x)\dd x\to\infty$. Derivatives and products with polynomials of tempered distributions are tempered.

\begin{definition}
The Fourier transform of $T\in\mathcal S'(\R^d)$ is $\langle\hat T,\varphi\rangle=\langle T,\hat\varphi\rangle$.
\end{definition}

For $T=T_f$ with $f\in L^1$ this agrees with $\hat f$, by the multiplication formula (Proposition~\ref{prop:ftl1}(e)); for $f\in L^2$ it agrees with Plancherel's $\mathcal Ff$, by continuity, as the multiplication formula extends to $L^2$.

\begin{theorem}\label{thm:ftsprime}
$\mathcal F$ is a bijection of $\mathcal S'(\R^d)$, with $\mathcal F^4=I$, and $\widehat{\partial^\beta T}=(2\pi i\xi)^\beta\hat T$, $\widehat{(-2\pi ix)^\alpha T}=\partial^\alpha\hat T$. Moreover $\hat\delta=1$ and $\hat1=\delta$.
\end{theorem}
\begin{proof}
From the proof of Proposition~\ref{prop:schwartz},
\[\sup|\xi^\beta\partial^\alpha\hat\varphi|\le(2\pi)^{|\alpha|-|\beta|}\|\partial^\beta(x^\alpha\varphi)\|_1,\]
and $\|\psi\|_1\le C\sup|(1+|x|^2)^d\psi|$. Expanding $\partial^\beta(x^\alpha\varphi)$ by Leibniz, $p_N(\hat\varphi)\le C_Np_{N+2d}(\varphi)$. So
$|\langle\hat T,\varphi\rangle|\le Cp_N(\hat\varphi)\le C'p_{N+2d}(\varphi)$, and $\hat T\in\mathcal S'$. $\mathcal F^4=I$ on $\mathcal S$ (Proposition~\ref{prop:schwartz}), hence on $\mathcal S'$, and $\mathcal F$ is a bijection. The formulas follow from those on
$\mathcal S$: for instance $\langle\widehat{\partial_jT},\varphi\rangle=-\langle T,\partial_j\hat\varphi\rangle=\langle T,\widehat{2\pi ix_j\varphi}\rangle=\langle(2\pi i\xi_j)\hat T,\varphi\rangle$. Finally $\langle\hat\delta,\varphi\rangle=\hat\varphi(0)=\int\varphi$, so
$\hat\delta=1$, and $\hat1=\mathcal F^2\delta=\delta$ since $\delta$ is even.
\end{proof}

\begin{beauty}
The identity $\hat1=\delta$ is the meaning of the formula $\int e^{-2\pi ix\xi}\dd x=\delta(\xi)$ that physicists write without blinking: the integral diverges, but as a distribution it is exact. It says that the
exponentials $e_\xi(x)=e^{2\pi ix\xi}$ are ``orthonormal'' in the continuous sense, $\langle e_\xi,e_\eta\rangle=\delta(\xi-\eta)$, which is the precise form of the ``as if'' of Plancherel's theorem in Chapter~5.
\end{beauty}

\section{Fundamental solutions}

A \emph{fundamental solution} of a linear differential operator $P(\partial)$ with constant coefficients is a distribution $E$ with $P(\partial)E=\delta$. Then $u=E*f$ solves $P(\partial)u=f$, at least for $f\in\mathcal D$, since
$P(\partial)(E*f)=(P(\partial)E)*f=\delta*f=f$ by Proposition~\ref{prop:smoothing}. On $\R$, $H$ is a fundamental solution of $\frac{\dd}{\dd x}$, and $u=H*f$ is the primitive $\int_{-\infty}^xf$. The Malgrange--Ehrenpreis theorem
(1954--1955) says that every non-zero operator with constant coefficients has one. Here is the most important.

\begin{theorem}[Newton's potential]\label{thm:newton}
In $\R^3$, $E(x)=-\frac1{4\pi|x|}$ is locally integrable and $\Delta E=\delta$.
\end{theorem}
\begin{proof}
In spherical coordinates $\int_{|x|\le R}\frac{\dd x}{|x|}=\int_0^R4\pi r^2\frac{\dd r}r<\infty$, so $E\in L^1_{\mathrm{loc}}$. For $\varepsilon>0$ let $E_\varepsilon(x)=-\frac1{4\pi}(|x|^2+\varepsilon^2)^{-1/2}$, which is smooth. For a smooth radial function $g(|x|)$
on $\R^3$, the chain rule gives $\Delta[g(|x|)]=g''(r)+\frac2rg'(r)$ for $r=|x|\ne0$. With $g(r)=-\frac1{4\pi}(r^2+\varepsilon^2)^{-1/2}$,
\[g'=\frac{r}{4\pi(r^2+\varepsilon^2)^{3/2}},\qquad g''+\frac2rg'=\frac{3(r^2+\varepsilon^2)-3r^2}{4\pi(r^2+\varepsilon^2)^{5/2}},\]
so $\Delta E_\varepsilon=\psi_\varepsilon$ with $\psi_\varepsilon(x)=\frac{3\varepsilon^2}{4\pi(|x|^2+\varepsilon^2)^{5/2}}=\varepsilon^{-3}\psi_1(\frac x\varepsilon)$ (by continuity, also at $0$). And $\int\psi_1=3\int_0^\infty\frac{r^2\dd r}{(1+r^2)^{5/2}}=1$ (substitute
$r=\tan u$: $\int_0^{\pi/2}\sin^2u\cos u\dd u=\frac13$). So $\psi_\varepsilon$ is an approximate identity. For $\varphi\in\mathcal D(\R^3)$, integrating by parts twice (no boundary terms),
\[\langle\Delta E,\varphi\rangle=\int E\Delta\varphi=\lim_{\varepsilon\to0}\int E_\varepsilon\Delta\varphi=\lim_{\varepsilon\to0}\int\psi_\varepsilon\varphi=\varphi(0),\]
by dominated convergence ($|E_\varepsilon|\le|E|$) and Tome~6, Theorem~9.10(b).
\end{proof}

So the solution of Poisson's equation $\Delta u=f$ is $u=E*f=-\frac1{4\pi}\int\frac{f(y)}{|x-y|}\dd y$: the gravitational potential of a mass density $f$, exactly as Newton and Laplace computed it, now with $\Delta u=f$ holding in
the sense of distributions for every $f\in\mathcal D$. The fundamental solution of the heat operator is the heat kernel of Chapter~5 (Exercise~\ref{ex:heatfund}).

\section{Sobolev spaces in one dimension}

\begin{definition}
For $s\ge0$, $H^s(\R)=\{f\in L^2(\R):(1+\xi^2)^{s/2}\hat f\in L^2\}$, with $\|f\|_{H^s}=\|(1+\xi^2)^{s/2}\hat f\|_2$. For an open interval $I$ and $1\le p\le\infty$,
$W^{1,p}(I)=\{u\in L^p(I):u'\in L^p(I)\}$, the derivative taken in $\mathcal D'(I)$, with $\|u\|_{W^{1,p}}=\|u\|_p+\|u'\|_p$; and $H^1(I)=W^{1,2}(I)$ with $\|u\|_{H^1}^2=\|u\|_2^2+\|u'\|_2^2$.
\end{definition}

By Theorem~\ref{thm:ftsprime} and Plancherel, $f\in L^2(\R)$ is in $H^1(\R)$ if and only if $f'\in L^2$, and then $\|f\|_{H^1}^2=\|f\|_2^2+\frac1{4\pi^2}\|f'\|_2^2$: two equivalent definitions, one counting derivatives, one
measuring the decay of $\hat f$. The Fourier one allows fractional $s$. $H^s(\R)$ is a Hilbert space, isometric to a weighted $L^2$ space, and $H^1(I)$ is complete by the same argument as for $L^2$ (if $u_n$ and $u_n'$
converge in $L^2$ to $u$ and $v$, then $u'=v$ by testing on $\varphi$).

\begin{theorem}[Sobolev embedding]\label{thm:sobolevR}
If $s>\frac12$, every $f\in H^s(\R)$ is almost everywhere equal to a continuous function tending to $0$ at infinity, and $\|f\|_\infty\le C_s\|f\|_{H^s}$ with $C_s^2=\int(1+\xi^2)^{-s}\dd\xi$.
\end{theorem}
\begin{proof}
By Cauchy--Schwarz, $\int|\hat f|=\int(1+\xi^2)^{-s/2}\cdot(1+\xi^2)^{s/2}|\hat f|\le C_s\|f\|_{H^s}$, finite since $2s>1$. So $\hat f\in L^1\cap L^2$, and by Theorems~\ref{thm:ftinv} and~\ref{thm:plancherel}, $f$ is almost everywhere equal to
$x\mapsto\int\hat f(\xi)e^{2\pi ix\xi}\dd\xi$, which is continuous, bounded by $\|\hat f\|_1$, and tends to $0$ at infinity by Riemann--Lebesgue.
\end{proof}

At $s=\frac12$ the embedding fails (Exercise~\ref{ex:loglog}). In dimension $d$ the threshold is $s>\frac d2$, with the same proof.

\begin{theorem}\label{thm:w11}
Let $I=]a,b[$ be bounded and $u\in W^{1,1}(I)$. Then $u$ has a representative $\tilde u$, continuous on $[a,b]$, with $\tilde u(y)-\tilde u(x)=\int_x^yu'$ for all $x,y$. If $u\in H^1(I)$, then
\[|\tilde u(y)-\tilde u(x)|\le|y-x|^{1/2}\|u'\|_2,\qquad\|\tilde u\|_\infty\le|I|^{-1/2}\|u\|_2+|I|^{1/2}\|u'\|_2.\]
\end{theorem}
\begin{proof}
Let $v(x)=\int_a^xu'$, continuous on $[a,b]$ since $u'\in L^1$. For $\varphi\in\mathcal D(I)$, Fubini gives
\[\int_av\varphi'=\int_a^b\int_a^xu'(t)\dd t\,\varphi'(x)\dd x=\int_a^bu'(t)\int_t^b\varphi'(x)\dd x\dd t=-\int_a^bu'\varphi,\]
so $v'=u'$ in $\mathcal D'(I)$. By Theorem~\ref{thm:zeroderiv}, $u-v$ is a constant $c$ almost everywhere, and $\tilde u=v+c$ works. For $u\in H^1$, Cauchy--Schwarz gives $|\int_x^yu'|\le|y-x|^{1/2}\|u'\|_2$. Since the mean of
$|\tilde u|^2$ over $I$ is $\frac{\|u\|_2^2}{|I|}$, there is $x_0$ with $|\tilde u(x_0)|\le|I|^{-1/2}\|u\|_2$, and $|\tilde u(x)|\le|\tilde u(x_0)|+|I|^{1/2}\|u'\|_2$.
\end{proof}

From now on we identify $u\in W^{1,1}(I)$ with its continuous representative. So functions of $H^1$ on an interval are H\"older continuous of exponent $\frac12$, and point values $u(x)$, in particular boundary values, make
sense and depend continuously on $u$ in $H^1$. Define $H^1_0(I)=\{u\in H^1(I):u(a)=u(b)=0\}$, a closed subspace.

\begin{theorem}[Rellich]\label{thm:rellich}
Let $I$ be a bounded interval. Every bounded sequence in $H^1(I)$ has a subsequence that converges uniformly on $\bar I$, hence in $L^2(I)$. So the inclusion $H^1(I)\to L^2(I)$ is a compact operator.
\end{theorem}
\begin{proof}
If $\|u_n\|_{H^1}\le M$, Theorem~\ref{thm:w11} shows that the $u_n$ are uniformly bounded and equicontinuous (equi-H\"older) on $\bar I$. By Arzel\`a--Ascoli (Tome~6, Theorem~4.15), a subsequence converges uniformly, and
uniform convergence on a bounded interval implies convergence in $L^2$.
\end{proof}

\begin{think}
The derivative of $|x|$ is the sign function, and the derivative of the sign is $2\delta$. So $|x|$ is in $H^1(]-1,1[)$, but the sign function is not. Is $\sqrt{|x|}$ in $H^1(]-1,1[)$? Is $|x|^{3/4}$?
\end{think}

\section{Weak solutions and the Lax--Milgram theorem}

Let $I=]0,1[$ and consider the boundary value problem
\begin{equation}\label{eq:bvp}
-(pu')'+qu=f\ \text{ on }I,\qquad u(0)=u(1)=0,
\end{equation}
with $p,q\in L^\infty(I)$, $p\ge p_0>0$, $q\ge0$, and $f\in L^2(I)$. The coefficients may jump, as for a rod made of two materials, so classical solutions may not exist. Multiply by $v\in H^1_0(I)$ and integrate by
parts, pretending everything is smooth: $\int pu'v'+\int quv=\int fv$. This makes sense for $u\in H^1_0$.

\begin{definition}
A \emph{weak solution} of \eqref{eq:bvp} is $u\in H^1_0(I)$ with $a(u,v)=\int fv$ for all $v\in H^1_0(I)$, where $a(u,v)=\int_0^1(pu'v'+quv)$.
\end{definition}

We work with real-valued functions in this section; everything holds for complex ones with sesquilinear forms.

\begin{theorem}[Lax--Milgram, 1954]\label{thm:laxmilgram}
Let $V$ be a real Hilbert space, $a:V\times V\to\R$ bilinear, \emph{bounded} ($|a(u,v)|\le M\|u\|\|v\|$) and \emph{coercive} ($a(u,u)\ge\alpha\|u\|^2$ with $\alpha>0$). For every bounded linear form $\ell$ on $V$ there is a unique
$u\in V$ with $a(u,v)=\ell(v)$ for all $v\in V$, and $\|u\|\le\frac{\|\ell\|}\alpha$.
\end{theorem}
\begin{proof}
For fixed $u$, $v\mapsto a(u,v)$ is a bounded linear form, so by the Riesz representation theorem (Theorem~\ref{thm:rieszrep}) there is a unique $Au\in V$ with $a(u,v)=\langle Au,v\rangle$; $A$ is linear and $\|Au\|\le M\|u\|$. Coercivity gives
$\alpha\|u\|^2\le\langle Au,u\rangle\le\|Au\|\|u\|$, so $\|Au\|\ge\alpha\|u\|$: $A$ is injective with closed range (as in Proposition~\ref{prop:saspec}). If $w\perp\operatorname{ran}A$, then $0=\langle Aw,w\rangle\ge\alpha\|w\|^2$, so the range is dense,
hence all of $V$. Let $y$ represent $\ell$ (Theorem~\ref{thm:rieszrep}) and $u=A^{-1}y$: then $a(u,v)=\langle y,v\rangle=\ell(v)$. Uniqueness follows from injectivity, and $\alpha\|u\|^2\le a(u,u)=\ell(u)\le\|\ell\|\|u\|$.
\end{proof}

When $a$ is symmetric, it is an inner product equivalent to the original one, and the theorem is just Riesz's; its power is for non-symmetric problems (Exercise~\ref{ex:convdiff}).

\begin{theorem}\label{thm:weaksol}
Problem \eqref{eq:bvp} has a unique weak solution $u$. It satisfies $\|u\|_{H^1}\le C\|f\|_2$; the function $pu'$ belongs to $H^1(I)$ and $-(pu')'+qu=f$ almost everywhere. If $p=1$ and $q,f$ are continuous, then
$u\in C^2([0,1])$ is a classical solution.
\end{theorem}
\begin{proof}
On $H^1_0(I)$, $|u(x)|=|\int_0^xu'|\le\|u'\|_2$, so $\|u\|_2\le\|u'\|_2$ (a Poincar\'e inequality; Exercise~\ref{ex:poincare} has the best constant). Hence $a(u,u)\ge p_0\|u'\|_2^2\ge\frac{p_0}2\|u\|_{H^1}^2$, and
$|a(u,v)|\le(\|p\|_\infty+\|q\|_\infty)\|u\|_{H^1}\|v\|_{H^1}$. The form $\ell(v)=\int fv$ has $\|\ell\|\le\|f\|_2$. Lax--Milgram gives existence, uniqueness and the bound with $C=\frac2{p_0}$.

Taking $v=\varphi\in\mathcal D(I)$ in the weak formulation, $\int pu'\varphi'=\int(f-qu)\varphi$, which says that the distributional derivative of $pu'\in L^2$ is $qu-f\in L^2$. So $pu'\in H^1(I)$, continuous, and
$(pu')'=qu-f$ almost everywhere. If $p=1$ and $q,f$ are continuous, then $u'$ is continuous with $(u')'=qu-f$ continuous (Theorem~\ref{thm:zeroderiv}), so $u\in C^2$.
\end{proof}

\begin{remark}
This is the modern route to Theorem~\ref{thm:sturm}: existence first, in a Hilbert space, with almost no hypotheses; regularity afterwards. The solution operator $f\mapsto u$ is bounded from $L^2$ to $H^1_0$, hence compact from $L^2$ to $L^2$
by Rellich, and self-adjoint (by the symmetry of $a$). The spectral theorem then gives a Hilbert basis of eigenfunctions of $-(pu')'+qu$ even for discontinuous coefficients. In dimension $d$, the same scheme solves
the Dirichlet problem for elliptic equations on bounded domains, with $H^1_0(\Omega)$ defined as the closure of $\mathcal D(\Omega)$.
\end{remark}

\emph{Galerkin's method.} Choose a finite-dimensional subspace $V_h\subset V$ and solve $a(u_h,v)=\ell(v)$ for $v\in V_h$, a linear system. By Lax--Milgram in $V_h$ it has a unique solution.

\begin{proposition}[C\'ea's lemma]\label{prop:cea}
In the setting of Theorem~\ref{thm:laxmilgram}, $\|u-u_h\|\le\frac M\alpha\inf_{v\in V_h}\|u-v\|$.
\end{proposition}
\begin{proof}
Subtracting, $a(u-u_h,w)=0$ for $w\in V_h$ (\emph{Galerkin orthogonality}). For $v\in V_h$, $\alpha\|u-u_h\|^2\le a(u-u_h,u-u_h)=a(u-u_h,u-v)\le M\|u-u_h\|\|u-v\|$.
\end{proof}

So the Galerkin solution is as good as the best approximation from $V_h$, up to the factor $\frac M\alpha$. The \emph{finite element method} takes for $V_h$ the continuous functions that are affine on each interval of a
mesh $0=x_0<x_1<\cdots<x_n=1$ and vanish at $0$ and $1$, with the ``hat functions'' as basis (Figure~\ref{fig:hats}). The matrix of $a$ in this basis is tridiagonal, and if $u\in H^2$ the error in $H^1$ is $O(h)$ for a
mesh of size $h$ (a numerical calculation measures it). Galerkin proposed the method in 1915 for problems of elasticity; Courant suggested piecewise linear trial functions in 1943, and engineers rediscovered it in the 1950s
to design aircraft.

\begin{figure}[ht]
\centering
\fig{ch06-hats}
\caption{The hat functions $\phi_1,\dots,\phi_5$ on a uniform mesh of $[0,1]$ (above), and a function of $V_h$, the combination $\sum_iu_i\phi_i$ with $u_i$ its values at the nodes (below). Each $\phi_i$ meets only its
neighbours, so the matrix $a(\phi_i,\phi_j)$ is tridiagonal.}
\label{fig:hats}
\end{figure}

\section{Exercises}

\begin{exercise}[$\star$]\label{ex:derivabs}
Compute, in $\mathcal D'(\R)$, the first and second derivatives of $|x|$, of $x_+=\max(x,0)$ and of $\sin|x|$. Show that $x\delta=0$, $x\delta'=-\delta$ and $x\operatorname{p.v.}\frac1x=1$. Deduce that the product of distributions cannot be
associative and commutative in general: compute $(\delta\cdot x)\cdot\operatorname{p.v.}\frac1x$ and $\delta\cdot(x\cdot\operatorname{p.v.}\frac1x)$ as if it were.
\end{exercise}

\begin{exercise}[$\star\star$]\label{ex:pvorder}
\emph{The principal value.} (a) Show that $\operatorname{p.v.}\frac1x$ is not of order $0$: find test functions $\varphi_n$ with $\|\varphi_n\|_\infty\le1$, supported in $[-1,1]$, and $\langle\operatorname{p.v.}\frac1x,\varphi_n\rangle\to\infty$. (b) Show that
$\widehat{\operatorname{p.v.}\frac1x}=-i\pi\operatorname{sgn}\xi$. (The operator $f\mapsto\mathcal F^{-1}(-i\operatorname{sgn}\xi\,\hat f)$, which is $\frac1\pi\operatorname{p.v.}\frac1x*f$, is the \emph{Hilbert transform}.)
\end{exercise}

\begin{exercise}[$\star\star$]\label{ex:loglog}
\emph{The failure at $s=\frac12$.} Let $g(\xi)=\frac1{(1+|\xi|)\log(2+|\xi|)}$ and $f=\mathcal F^{-1}g$. Show that $f\in H^{1/2}(\R)$ but that $f$ is not essentially bounded.
\end{exercise}

\begin{exercise}[$\star$]\label{ex:poincare}
\emph{Poincar\'e's inequality, sharp.} Show that $\|u\|_2\le\frac1\pi\|u'\|_2$ for $u\in H^1_0(]0,1[)$, with equality for $u=\sin\pi x$. Relate the constant to Theorem~\ref{thm:sturm}.
\end{exercise}

\begin{exercise}[$\star$]\label{ex:gagliardo}
\emph{An interpolation inequality.} Show that $\|u\|_\infty^2\le\|u\|_2\|u'\|_2$ for real $u\in H^1(\R)$, and that equality holds for $u=e^{-|x|}$. (Approximate by test functions first.)
\end{exercise}

\begin{exercise}[$\star\star$]\label{ex:convdiff}
\emph{Convection and diffusion.} Let $b\in\R$ and $f\in L^2(]0,1[)$. Show that $-u''+bu'+u=f$, $u(0)=u(1)=0$, has a unique weak solution in $H^1_0$, although the bilinear form is not symmetric. For $f=1$, solve it explicitly and describe
the solution when $b$ is large: where is the ``boundary layer''?
\end{exercise}

\begin{exercise}[$\star\star$]\label{ex:heatfund}
\emph{The heat kernel as a fundamental solution.} On $\R^2$ with coordinates $(t,x)$, let $K(t,x)=\frac{\mathbf 1_{t>0}}{\sqrt{4\pi t}}e^{-x^2/(4t)}$. Show that $K\in L^1_{\mathrm{loc}}(\R^2)$ and $(\partial_t-\partial_x^2)K=\delta_{(0,0)}$.
\end{exercise}

\section{Hints}

\begin{hint}{ex:derivabs}
$|x|'=\operatorname{sgn}x$ and $\operatorname{sgn}'=2\delta$. For products, test against $\varphi$ and move $x$ onto it.
\end{hint}
\begin{hint}{ex:pvorder}
(a) Take $\varphi_n$ close to $\operatorname{sgn}x$ on $[\frac1n,1]$ and $[-1,-\frac1n]$. (b) Approximate by the truncations $\mathbf 1_{\varepsilon<|x|<1/\varepsilon}\frac1x$ and use $\int_0^\infty\frac{\sin t}t\dd t=\frac\pi2$.
\end{hint}
\begin{hint}{ex:loglog}
$\int\frac{\dd\xi}{(1+|\xi|)\log^2(2+|\xi|)}<\infty$ but $\int g=\infty$. If $|f|\le M$, look at $f*G_\varepsilon(0)$.
\end{hint}
\begin{hint}{ex:poincare}
Expand $u$ in the sine basis of Lemma~\ref{lem:sines}, and $u'$ in the cosines $\sqrt2\cos(n\pi x)$.
\end{hint}
\begin{hint}{ex:gagliardo}
$u(x)^2=\int_{-\infty}^x(u^2)'=-\int_x^\infty(u^2)'$.
\end{hint}
\begin{hint}{ex:convdiff}
$\int_0^1u'u=\frac12[u^2]_0^1=0$ on $H^1_0$, so $a(u,u)=\|u\|_{H^1}^2$.
\end{hint}
\begin{hint}{ex:heatfund}
\[\langle K,-\partial_t\varphi-\partial_x^2\varphi\rangle=\lim_{\varepsilon\to0}\int_\varepsilon^\infty\!\!\int K(-\partial_t\varphi-\partial_x^2\varphi);\]
integrate by parts in $t$ and in $x$.
\end{hint}

\section{Solutions}

\begin{solution}{ex:derivabs}
\idea A distributional derivative is the classical derivative where it exists, plus a Dirac mass for every jump. Corners produce jumps in the first derivative, hence masses in the second.
\writeup $|x|$ is continuous, with classical derivative $\operatorname{sgn}x$ off $0$; integrating by parts on each half-line, $|x|'=\operatorname{sgn}x$, and $\operatorname{sgn}$ jumps by $2$ at $0$, so $|x|''=2\delta$. Likewise $x_+'=H$ and
$x_+''=\delta$. For $\sin|x|$: the first derivative is $\operatorname{sgn}(x)\cos x$, which jumps by $2$ at $0$, so the second is $-\sin|x|+2\delta$. Next, $\langle x\delta,\varphi\rangle=(x\varphi)(0)=0$;
$\langle x\delta',\varphi\rangle=-(x\varphi)'(0)=-\varphi(0)$; and $\langle x\operatorname{p.v.}\frac1x,\varphi\rangle=\lim\int_{|x|>\varepsilon}\varphi=\int\varphi$. If a product of distributions extended these rules associatively and
commutatively, then $(\delta\cdot x)\cdot\operatorname{p.v.}\frac1x=0\cdot\operatorname{p.v.}\frac1x=0$, while $\delta\cdot(x\cdot\operatorname{p.v.}\frac1x)=\delta\cdot1=\delta$. Schwartz showed in 1954 that no multiplication on all distributions
extends that of continuous functions and satisfies the usual algebraic rules; nonlinear problems need other tools.
\end{solution}

\begin{solution}{ex:pvorder}
\idea The principal value sees the derivative of $\varphi$ at $0$, through the cancellation between $x$ and $-x$. An odd function with a steep jump at $0$ defeats any bound by $\|\varphi\|_\infty$ alone.
\writeup (a) Let $\varphi_n$ be odd, with $0\le\varphi_n\le1$ on $[0,1]$, $\varphi_n=1$ on $[\frac1n,\frac12]$, and support in $[-1,1]$. Then $\langle\operatorname{p.v.}\frac1x,\varphi_n\rangle=2\int_0^1\frac{\varphi_n(x)}x\dd x\ge2\log\frac n2\to\infty$, while
$\|\varphi_n\|_\infty=1$.

(b) Let $T_\varepsilon=\mathbf 1_{\varepsilon<|x|<1/\varepsilon}\frac1x\in L^1$. For $\varphi\in\mathcal S$, $\langle T_\varepsilon,\varphi\rangle=\int_\varepsilon^{1/\varepsilon}\frac{\varphi(x)-\varphi(-x)}x\dd x\to\langle\operatorname{p.v.}\frac1x,\varphi\rangle$ (dominated convergence, the
integrand being bounded by $2\|\varphi'\|_\infty$ near $0$ and by $\frac{2\|\varphi\|_\infty}x\cdot\frac{C}{1+x}$ far away). By oddness,
\[\hat T_\varepsilon(\xi)=-2i\int_\varepsilon^{1/\varepsilon}\frac{\sin(2\pi x\xi)}x\dd x=-2i\operatorname{sgn}(\xi)\int_{2\pi\varepsilon|\xi|}^{2\pi|\xi|/\varepsilon}\frac{\sin t}t\dd t.\]
The integrals $\int_a^b\frac{\sin t}t\dd t$ are bounded uniformly in $0\le a<b$ and tend to $\frac\pi2$ as $a\to0$, $b\to\infty$ (Tome~4, Chapter~4). So $\hat T_\varepsilon(\xi)\to-i\pi\operatorname{sgn}\xi$ boundedly, and for
$\varphi\in\mathcal S$, $\langle\widehat{\operatorname{p.v.}\frac1x},\varphi\rangle=\lim\langle T_\varepsilon,\hat\varphi\rangle=\lim\int\hat T_\varepsilon\varphi=\int-i\pi\operatorname{sgn}(\xi)\varphi(\xi)\dd\xi$, using the multiplication formula and dominated convergence.
\end{solution}

\begin{solution}{ex:loglog}
\idea The embedding at $s>\frac12$ came from $\hat f\in L^1$. At $s=\frac12$, one can make $\hat f$ just fail to be integrable while $|\xi|^{1/2}\hat f$ is still square-integrable, the difference being one power of a
logarithm.
\writeup $(1+\xi^2)^{1/2}g(\xi)^2\le\frac{\sqrt2}{(1+|\xi|)\log^2(2+|\xi|)}$, which is integrable (substitute $t=\log(2+|\xi|)$), so $f\in H^{1/2}$; also $g\in L^2$, so $f\in L^2$ is well defined. But $\int g=\infty$, since
$\int^\infty\frac{\dd\xi}{\xi\log\xi}$ diverges. Suppose $|f|\le M$ almost everywhere. With $G_\varepsilon$ as in Chapter~5, $f*G_\varepsilon(0)=\int g(\xi)e^{-\pi\varepsilon^2\xi^2}\dd\xi$ (Plancherel's inner product with the even function $G_\varepsilon$,
whose transform is $e^{-\pi\varepsilon^2\xi^2}$), and $|f*G_\varepsilon(0)|\le M$. But the right side tends to $\int g=\infty$ by monotone convergence: contradiction.
\end{solution}

\begin{solution}{ex:poincare}
\idea On $H^1_0$, differentiation is diagonal in the sine basis, multiplying the $n$-th coefficient by $n\pi$. The worst ratio $\frac{\|u\|}{\|u'\|}$ is at the slowest sine.
\writeup Let $c_n=\langle u,s_n\rangle$ with $s_n=\sqrt2\sin(n\pi x)$, and $\gamma_n(x)=\sqrt2\cos(n\pi x)$. Integrating by parts (valid for the absolutely continuous $u$, with $u(0)=u(1)=0$),
$\langle u',\gamma_n\rangle=-\int u\gamma_n'=n\pi c_n$. The $\gamma_n$, $n\ge1$, are orthonormal, so by Bessel and Parseval (Lemma~\ref{lem:sines}),
\[\|u'\|_2^2\ge\sum_{n\ge1}n^2\pi^2|c_n|^2\ge\pi^2\sum_n|c_n|^2=\pi^2\|u\|_2^2.\]
Equality holds for $u=s_1$. The constant $\pi^2$ is the first eigenvalue of $-u''$ with Dirichlet conditions (Theorem~\ref{thm:sturm} with $q=0$): the best Poincar\'e constant is always the bottom of the spectrum.
\end{solution}

\begin{solution}{ex:gagliardo}
\idea The value at a point is controlled by integrating the derivative of $u^2$ from either side; averaging the two sides gives the factor that makes the inequality sharp.
\writeup For $u\in\mathcal D(\R)$ real, $u(x)^2=\int_{-\infty}^x2uu'=-\int_x^\infty2uu'$, so $2u(x)^2\le\int2|u||u'|$, and Cauchy--Schwarz gives $u(x)^2\le\|u\|_2\|u'\|_2$. Test functions are dense in $H^1(\R)$ (truncate with
$\chi(\frac xn)$ and mollify), and both sides are continuous for the $H^1$ norm, the left one by Theorem~\ref{thm:sobolevR} with $s=1$. For $u=e^{-|x|}$, $\|u\|_\infty=1$, $\|u\|_2^2=\int e^{-2|x|}=1$ and $\|u'\|_2^2=1$: equality.
\end{solution}

\begin{solution}{ex:convdiff}
\idea Coercivity does not need symmetry. On functions vanishing at both ends, the convection term $bu'$ does no work against $u$, so the energy estimate is the same as for $b=0$; Lax--Milgram then replaces Riesz.
\writeup Let $a(u,v)=\int(u'v'+bu'v+uv)$ on $V=H^1_0$. It is bilinear and bounded by $(1+|b|)\|u\|_{H^1}\|v\|_{H^1}$. Since $uu'=\frac12(u^2)'$ and $u$ vanishes at both ends (Theorem~\ref{thm:w11}, applied to $u^2\in W^{1,1}$),
$\int bu'u=0$ and $a(u,u)=\|u\|_{H^1}^2$. Lax--Milgram gives a unique weak solution, and as in Theorem~\ref{thm:weaksol}, $u'\in H^1$ with $-u''+bu'+u=f$ almost everywhere. For $f=1$: $u=1+Ae^{r_1x}+Be^{r_2x}$ with
$r_{1,2}=\frac{b\mp\sqrt{b^2+4}}2$ and $A,B$ fixed by $u(0)=u(1)=0$. For $b\gg1$, $r_1\approx-\frac1b$ is small and $r_2\approx b$ is large. Away from $x=1$ the term $e^{r_2x}$ is negligible
(its coefficient is of size $e^{-b}$), and $u$ follows the \emph{reduced} first-order equation $bu'+u=1$ with $u(0)=0$, that is $u\approx1-e^{-x/b}\approx\frac xb$: the diffusion $-u''$ has dropped out. A first-order
equation can only honour one boundary condition, and the flow carries information from left to right, so it keeps the inflow condition at $x=0$. The outflow condition $u(1)=0$ is restored by the fast exponential
$e^{b(x-1)}$, which lives in a layer of width about $\frac1b$ at $x=1$: that is the boundary layer. (With $f=b$ instead, $u\approx x$ and the drop from $1$ to $0$ inside the layer is dramatic.) Numerically,
such layers make naive finite elements oscillate unless the mesh resolves them, a famous difficulty of convection-dominated problems.
\end{solution}

\begin{solution}{ex:heatfund}
\idea The heat kernel is the temperature after a unit of heat is dropped at the origin at time $0$. Testing against $\varphi$ and cutting out the instant $t=0$, the equation $(\partial_t-\partial_x^2)K=0$ kills the inside, and what survives is the
temperature profile at time $\varepsilon$, which is an approximate identity.
\writeup \emph{Local integrability.} $K\ge0$ and $\int K(t,x)\dd x=1$ for every $t>0$ (Gaussian integral), so $\int_0^T\int K=T$: $K\in L^1_{\mathrm{loc}}(\R^2)$ (even integrable on every strip $t\le T$).

\emph{The equation.} We must show $\langle K,-\partial_t\varphi-\partial_x^2\varphi\rangle=\varphi(0,0)$ for $\varphi\in\mathcal D(\R^2)$. By dominated convergence, the left side is
$\lim_{\varepsilon\to0}I_\varepsilon$ with $I_\varepsilon=\int_\varepsilon^\infty\int K(-\partial_t\varphi-\partial_x^2\varphi)\dd x\dd t$. On $t\ge\varepsilon$, $K$ is smooth and all its derivatives decay like Gaussians in $x$, so we may
integrate by parts: twice in $x$ (no boundary terms, $\varphi$ having compact support), and once in $t$, with a boundary term at $t=\varepsilon$:
\[I_\varepsilon=\int K(\varepsilon,x)\varphi(\varepsilon,x)\dd x+\int_\varepsilon^\infty\!\!\int(\partial_tK-\partial_x^2K)\varphi\dd x\dd t.\]
A direct computation gives $\partial_tK=\partial_x^2K=\bigl(\frac{x^2}{4t^2}-\frac1{2t}\bigr)K$ for $t>0$, so the second term vanishes. In the first, $K(\varepsilon,\cdot)$ is the Gaussian $G$ of variance $2\varepsilon$, an approximate identity
as $\varepsilon\to0$, and
\[\Bigl|\int K(\varepsilon,x)\varphi(\varepsilon,x)\dd x-\varphi(0,0)\Bigr|\le\int K(\varepsilon,x)\bigl|\varphi(\varepsilon,x)-\varphi(0,0)\bigr|\dd x.\]
Split at $|x|=\eta$: where $|x|\le\eta$ the integrand is at most $\sup_{|s|,|y|\le\max(\varepsilon,\eta)}|\varphi(s,y)-\varphi(0,0)|$, small by uniform continuity; where $|x|>\eta$ it is at most $2\|\varphi\|_\infty\int_{|x|>\eta}K(\varepsilon,x)\dd x$,
which tends to $0$ as $\varepsilon\to0$. So $I_\varepsilon\to\varphi(0,0)$, that is $(\partial_t-\partial_x^2)K=\delta$. The solution of $\partial_tu=\partial_x^2u$ with $u(0,\cdot)=g$ is then $K(t,\cdot)*g$, as Fourier promised in
Exercise~\ref{ex:heatft}.
\end{solution}

\section*{Sources and further reading}
\addcontentsline{toc}{section}{Sources and further reading}

Laurent Schwartz announced distributions in 1945 and gave the full theory in \emph{Th\'eorie des distributions} (Hermann, 2 vols., 1950--51), which won him the Fields Medal in 1950. His autobiography, \emph{Un math\'ematicien aux prises avec
le si\`ecle} (Odile Jacob, 1997; English translation AMS 2001), tells the story of the ``night of the distributions'' and much else. The idea had precursors: Dirac's $\delta$ in \emph{The Principles of Quantum Mechanics}
(Oxford, 1930), Heaviside's operational calculus, and above all S.~L.~Sobolev, ``M\'ethode nouvelle \`a r\'esoudre le probl\`eme de Cauchy pour les \'equations lin\'eaires hyperboliques normales'', \emph{Mat.\ Sbornik} 1 (1936),
39--72, where generalized functions and the spaces now named after him first appear. The Lax--Milgram theorem is from P.~D.~Lax and A.~N.~Milgram, ``Parabolic equations'', in \emph{Contributions to the Theory of Partial
Differential Equations}, Ann.\ Math.\ Studies 33 (Princeton, 1954), 167--190. The finite element method has a famous ancestor in R.~Courant, ``Variational methods for the solution of problems of equilibrium and vibrations'',
\emph{Bull.\ AMS} 49 (1943), 1--23. Heaviside's dealings with the Royal Society are told in P.~J.~Nahin, \emph{Oliver Heaviside: Sage in Solitude} (IEEE Press, 1988). The existence of a fundamental solution for every constant-coefficient operator is due to Malgrange and to Ehrenpreis (1954--55).

For a first course, H.~Brezis, \emph{Functional Analysis, Sobolev Spaces and Partial Differential Equations} (Springer, 2011), Chapters~8--9, is the model of this chapter's one-dimensional approach. The standard reference for distributions is
L.~H\"ormander, \emph{The Analysis of Linear Partial Differential Operators I} (Springer, 1983). For finite elements, S.~C.~Brenner and L.~R.~Scott, \emph{The Mathematical Theory of Finite Element Methods} (Springer, 3rd ed., 2008).
